% GENERATED FILE -- do not edit by hand.
% Numbering of the main paper, extracted from its main.aux so that
% references into the paper render as real numbers without needing the
% paper tree at build time. Regenerate whenever the paper is renumbered.
\makeatletter
\newlabel{M:fig:family-census-tree}{{10}{52}{Anatomy of the 400 scored, accepted source commits from the public competition through 18 July (\#883--1294). Each commit is assigned to one primary optimization family (level~1) and, where applicable, one change type (level~2), based on an audit of its \texttt {git show} diff rather than its often-unreliable subject or lever name. Blue nodes represent nine circuit-optimization families: six conventional compiler categories and three reversibility-specific categories---uncompute strategy, register/scratch allocation, and \emph {codec/encoding}---with distinct resource and correctness obligations. Amber nodes represent verifier search for a clean Fiat--Shamir nonce, which is not a circuit optimization, and neutral or mechanical changes. Counts measure how often each mechanism was the primary change, not its contribution to score reduction; accompanying nonce changes are attributed to the associated circuit optimization. Pre-competition development commits are excluded, including the representation/coordinate-change family documented only in earlier research (see \href {https://github.com/jieyilong/ecdsafail-supplemental-materials/blob/main/appendices/detailed_results_and_circuit_improvement_trajectory.pdf}{Supplemental Note A: Detailed results and circuit-improvement trajectory}), because no corresponding competition submission was accepted. Four accepted but unscored \texttt {DUMMY\_TOFFOLIS} bootstrap commits are also excluded. The complete per-commit ledger, including those four commits, is provided in \href {https://github.com/jieyilong/ecdsafail-supplemental-materials/blob/main/appendices/optimization_census_of_accepted_competition_source_commits.pdf}{Supplemental Note C: List of accepted solutions}}{figure.caption.43}{}}
\newlabel{M:fig:family-census-tree@cref}{{[figure][10][]10}{[1][51][]52}{}{}{}}
\newlabel{M:sec:base-five-codec}{{5.3.2}{39}{Base-5 Codec for the Jump-2 Transcript}{subsubsection.5.3.2}{}}
\newlabel{M:sec:base-five-codec@cref}{{[subsubsection][2][5,3]5.3.2}{[1][39][]39}{}{}{}}
\newlabel{M:sec:consprop}{{5.3.6}{46}{Constant Propagation}{subsubsection.5.3.6}{}}
\newlabel{M:sec:consprop@cref}{{[subsubsection][6][5,3]5.3.6}{[1][46][]46}{}{}{}}
\newlabel{M:sec:cost}{{3.2}{17}{Circuit Model and Cost Metric}{subsection.3.2}{}}
\newlabel{M:sec:cost@cref}{{[subsection][2][3]3.2}{[1][17][]17}{}{}{}}
\newlabel{M:sec:dead-code-elimination}{{5.3.7}{46}{Dead Code Elimination}{subsubsection.5.3.7}{}}
\newlabel{M:sec:dead-code-elimination@cref}{{[subsubsection][7][5,3]5.3.7}{[1][46][]46}{}{}{}}
\newlabel{M:sec:ecc-quantum}{{2.2}{10}{Elliptic Curve Cryptography and Quantum Attacks}{subsection.2.2}{}}
\newlabel{M:sec:ecc-quantum@cref}{{[subsection][2][2]2.2}{[1][10][]10}{}{}{}}
\newlabel{M:sec:evaluator}{{3.3}{18}{Correctness, Verification, and Scoring}{subsection.3.3}{}}
\newlabel{M:sec:evaluator@cref}{{[subsection][3][3]3.3}{[1][18][]18}{}{}{}}
\newlabel{M:sec:jump-two-gcd}{{5.3.1}{37}{The Jump-2 Euclidean Algorithm}{subsubsection.5.3.1}{}}
\newlabel{M:sec:jump-two-gcd@cref}{{[subsubsection][1][5,3]5.3.1}{[1][37][]37}{}{}{}}
\newlabel{M:sec:limitations}{{6}{54}{Limitations, Safety, and Responsible Framing}{section.6}{}}
\newlabel{M:sec:limitations@cref}{{[section][6][]6}{[1][54][]54}{}{}{}}
\newlabel{M:sec:low-qubit-track}{{5.4.3}{50}{Low-Qubit and Pareto Exploration}{subsubsection.5.4.3}{}}
\newlabel{M:sec:low-qubit-track@cref}{{[subsubsection][3][5,4]5.4.3}{[1][50][]50}{}{}{}}
\newlabel{M:sec:optimization-technique-breakdown}{{5.4.4}{51}{Breakdown of Optimization Techniques Across the Circuit Optimization Trajectory}{subsubsection.5.4.4}{}}
\newlabel{M:sec:optimization-technique-breakdown@cref}{{[subsubsection][4][5,4]5.4.4}{[1][51][]51}{}{}{}}
\newlabel{M:sec:register-shared-eea}{{5.3.5}{43}{Register-Shared EEA and Lane-Lifetime Refinements}{subsubsection.5.3.5}{}}
\newlabel{M:sec:register-shared-eea@cref}{{[subsubsection][5][5,3]5.3.5}{[1][43][]43}{}{}{}}
\newlabel{M:sec:results-summary}{{5.1}{28}{Overall Summary of Results}{subsection.5.1}{}}
\newlabel{M:sec:results-summary@cref}{{[subsection][1][5]5.1}{[1][28][]28}{}{}{}}
\newlabel{M:sec:specialized-squaring}{{5.3.4}{42}{Karatsuba Squaring with Pseudo-Mersenne Modulo}{subsubsection.5.3.4}{}}
\newlabel{M:sec:specialized-squaring@cref}{{[subsubsection][4][5,3]5.3.4}{[1][42][]42}{}{}{}}
\newlabel{M:tab:admitted-pareto}{{2}{28}{Archive-conditioned Pareto set at the data cutoff. Adjacent differences are descriptive secants between distinct circuit families, not causal marginal effects. ``Admitted'' denotes an archived public clean-run report included under the manuscript's curated selection, not official promotion}{table.caption.30}{}}
\newlabel{M:tab:admitted-pareto@cref}{{[table][2][]2}{[1][28][]28}{}{}{}}
\newlabel{M:tab:external-pa-comparison}{{3}{31}{Reported resources, correctness evidence, and retry-adjusted spacetime-resource proxies for representative \texttt {secp256k1} point-addition circuits. Here $\mathrm {B}=10^9$. Under an independently rerunnable classical-success model, a circuit with empirical success probability $\hat p$ requires $1/\hat p$ executions on average to obtain one successful result. Accordingly, $Q\times T/\hat p$ provides a sensitivity-adjusted proxy when success probabilities and failure events are defined and measured comparably. Among the window-supporting circuits listed in the table, the {\challengename } windowed \texttt {8e9c9a2} construction has the lowest tabulated $Q\times T/\hat p$, at approximately $1.961\,\mathrm {B}$, while supporting the required single-call coherent window-indexed point-addition interface. To avoid penalizing external circuits for which no common-corpus estimate is available, rows marked ${}^{*}$ conservatively assume zero error, or $\hat p=1$; these are therefore optimistic estimates, and any nonzero error would increase the required work. The Google/Babbush entries remain upper bounds because their published $Q$ and $T$ values are themselves upper thresholds. FS denotes finite Fiat--Shamir support, whereas ``independent'' denotes a corpus selected independently of the circuit hash. Failure definitions and accounting conventions differ across sources, so the comparison remains contextual. Window support denotes a coherent window-indexed point-addition interface, not a complete implementation of Shor's algorithm. ${}^{\dagger }$The Schrottenloher values include the separately reported $w=16$ overhead of $16$ qubits and $3\cdot 2^{16}$ Toffoli gates; the corresponding arithmetic-core values are $(1{,}192,2.39\,\mathrm {M})$ and $(1{,}446,1.86\,\mathrm {M})$}{table.caption.31}{}}
\newlabel{M:tab:external-pa-comparison@cref}{{[table][3][]3}{[1][30][]31}{}{}{}}
\makeatother
